%------------------------------------------------------------------------------%
\documentclass[fleqn,utf8,aspectratio=169,12pt]{beamer}

\usepackage{calculo_varias_variaveis-1}

\title{Números Complexos}

%------------------------------------------------------------------------------%
\begin{document}

\addtitle

%------------------------------------------------------------------------------%
\section{Unidade Imaginária}

%------------------------------------------------------------------------------%
\begin{frame}{Motivação}
  Considere a equação
  \[ x^2 + 1 = 0 \]
  \pause
  Manipulando formalmente temos a solução
  \[ x = \sqrt{-1} \]
  \pause
  Porém, \(x^2 \geq 0\) para todo $x\in\R$

  \pause
  \bigskip
  Não existe \(\sqrt{-1}\) dentro dos números reais
\end{frame}

%------------------------------------------------------------------------------%
\begin{frame}{Unidade imaginária}
  \emph{Unidade imaginária}
  \[
    i^2 = -1
    \qquad\text{ou}\qquad
    i = \sqrt{-1}
  \]

  \pause
  A Engenharia Elétrica usa
  \bigskip
  \(j^2=-1\)
\end{frame}

%------------------------------------------------------------------------------%
\begin{frame}{Propriedades de Potências}
  \Large
  \[
    a^p a^q = a^{p+q}
  \]

  \pause
  \bigskip
  \[
    \frac{a^p}{a^q} = a^{p-q}
  \]

  \pause
  \bigskip
  \[
    \left(a^p\right)^q = a^{pq}
  \]
\end{frame}

%------------------------------------------------------------------------------%
\begin{frame}{Potências de $i$}
\large
\begin{columns}
\column{0.5\linewidth}
  {
  \setlength{\jot}{5pt}
  \begin{align*}
    \onslide<+->{i^1}
    \onslide<+->{ & = i } \\
    \onslide<+->{i^2}
    \onslide<+->{ & = -1 } \\
    \onslide<+->{i^3}
    \onslide<+->{& = i^2 \, i \;}
    \onslide<+->{ = (-1) i}
    \onslide<+->{ = -i}   \\
    \onslide<+->{i^4}
    \onslide<+->{ & = i^2 \, i^2}
    \onslide<+->{  = (-1)(-1)}
    \onslide<+->{ = 1}  \\
    \onslide<+->{i^5}
    \onslide<+->{ & = i^4 \, i \;}
    \onslide<+->{ = 1 i}
    \onslide<+->{ = i} \\
    \onslide<+->{i^6}
    \onslide<+->{ & = i^4 \, i^2}
    \onslide<+->{ = 1(-1)}
    \onslide<+->{ = -1}    \\
    \onslide<+->{i^7}
    \onslide<+->{ & = i^4 \, i^3}
    \onslide<+->{ = 1(-i)}
    \onslide<+->{ = -i}    \\
    \onslide<+->{i^8}
    \onslide<+->{ & = i^4 \, i^4}
    \onslide<+->{ = 1\times 1}
    \onslide<+->{ = 1}
  \end{align*}
  }
\column{0.5\linewidth}
\begin{align*}
  \onslide<+->{i^0 & = 1 } \\
  \onslide<+->{i^{-1}}
  \onslide<+->{& = \frac{1}{i}}
  \onslide<+->{  = \frac{i}{i^2}}
  \onslide<+->{  = \frac{i}{-1}}
  \onslide<+->{  = -i } \\
  \onslide<+->{i^{-2}}
  \onslide<+->{& = \frac{1}{i^2}}
  \onslide<+->{  = \frac{1}{-1}}
  \onslide<+->{  = -1 } \\
  \onslide<+->{i^{-3}}
  \onslide<+->{& = \frac{1}{i^3}}
  \onslide<+->{  = \frac{1}{-i}}
  \onslide<+->{  = \frac{i}{-i^2}}
  \onslide<+->{  = \frac{i}{1}}
  \onslide<+->{  = i } \\[7mm]
  \onslide<+->{i^{75}}
  \onslide<+->{& = ?}
\end{align*}
\end{columns}
\end{frame}

%------------------------------------------------------------------------------%
\begin{frame}{Potências de $i$}
\begin{columns}
\column{0.5\linewidth}
  \onslide<+->{Para uma potência $n$ inteira}

  \onslide<+->{
  \bigskip
  Dividimos $n$ por 4}
  \[
    \onslide<+->{
    \begin{array}{>{\;}c<{\;}>{\;}c<{\;}}
      n & \multicolumn{1}{|c}{4} \\ \cline{2-2}
      r & q
    \end{array}}
    \qquad
    \onslide<+->{
    \begin{array}{>{\;}c<{\;}>{\;}c<{\;}}
      75 & \multicolumn{1}{|c}{4} \\ \cline{2-2}
      3  & 18
    \end{array}}
  \]
  \onslide<+->{Assim}
  \begin{align*}
    \onslide<+->{n  & = 4q + r} \\
    \onslide<+->{75 & = 4\times 18 + 3}
  \end{align*}
\column{0.5\linewidth}
{\Large
  \begin{align*}
    \onslide<+->{i^n}
    \onslide<+->{& = i^{\,4q\,+\,r} } \\
    \onslide<+->{& = i^{\,4q} \; i^r } \\
    \onslide<+->{& = \left(i^4\right)^q \, i^r } \\
    \onslide<+->{& = \left(1  \right)^q \, i^r } \\
    \onslide<+->{& = i^r} \\[3mm]
    \onslide<+->{i^{75} & = i^3 = -i}
  \end{align*}
}
\end{columns}
\end{frame}

%------------------------------------------------------------------------------%
\begin{frame}{Raízes de Números Negativos}
  Seja \(a\in\R\), \(a>0\)
  \pause
  \[
    \sqrt{-a^2}
    \pause
    \;=\; \sqrt{(-1)\,a^2}
    \pause
    \;=\; i\sqrt{a^2}
    \pause
    \;=\; ai
  \]
\end{frame}

%------------------------------------------------------------------------------%
\begin{frame}{Cuidado com raízes}
\begin{columns}
\column{0.4\linewidth}
  Sabemos que
  \[
    \sqrt{a}\,\sqrt{b} = \sqrt{ab}
  \]

  \pause
  Por exemplo
  \[
    \sqrt{9}\,\sqrt{4}
    \pause
    = 3\times 2
    \pause
    = 6
  \]
  \[
    \pause
    \sqrt{9\times 4}
    \pause
    = \sqrt{36}
    \pause
    = 6
  \]
\column{0.6\linewidth}
  \pause
  \[
    \sqrt{-9}\,\sqrt{-4}
    \pause
    = 3i\times 2i
    \pause
    = 6i^2
    \pause
    = {\color<15->{red}-6}
  \]
  \[
    \pause
    \sqrt{(-9)(-4)}
    \pause
    = \sqrt{36}
    \pause
    = {\color<15->{red}6}
  \]

  \bigskip
  \onslide<16>{\emph{Só é verdade para $a$ e $b$ positivos!}}
\end{columns}
\end{frame}

%------------------------------------------------------------------------------%
%------------------------------------------------------------------------------%
\begin{question}
  Simplifique as expressões
  \Large
  \[
    a = i^7
  \]

  \[
    b = i^{57}
  \]

  \[
    c = \frac{i^5 + i^{17}}{i^3}
  \]
\end{question}

%------------------------------------------------------------------------------%
\begin{resolution}{Solução}
  \Large
  \[
    a
    \;=\; i^7
    \pause
    \;=\; i^4 \, i^3
    \pause
    \;=\; i^3
    \pause
    \;=\; i^2 i
    \pause
    \;=\; -i
  \]
\end{resolution}

%------------------------------------------------------------------------------%
\begin{resolution}{Solução}
  \Large
  \[
    b
    \;=\; i^{57}
    \pause
    \;=\; i^{4\times 14} \, i^1
    \pause
    \;=\; i
  \]
\end{resolution}

%------------------------------------------------------------------------------%
\begin{resolution}{Solução}
  \Large
  \[
    c
    \;=\; \frac{i^5 + i^{17}}{i^3}
    \pause
    \;=\; \frac{i^5}{i^3} + \frac{i^{17}}{i^3}
    \pause
    \;=\; i^{5-3} + i^{17-3}
  \]

  \pause
  \[
    \vphantom{c}
    \;=\; i^2 + i^{14}
    \pause
    \;=\; -1 + i^{\,4\times 3+2}
    \pause
    \;=\; -1 + i^2
  \]

  \pause
  \[
    \vphantom{c}
    \;=\; -1 -1
    \pause
    \;=\; -2
  \]
\end{resolution}

%------------------------------------------------------------------------------%


%------------------------------------------------------------------------------%
\section{Números Complexos}

%------------------------------------------------------------------------------%
\begin{frame}{Números complexos}
  Conjunto dos Números Complexos
  \[
    \mathbb{C} = \big\{\, z=a+bi \;\big|\; a, b \in \R \,\big\}
  \]
  \pause
  Parte real de $z$
  \[
    \Re(z) = a
  \]
  \pause
  Parte imaginária de $z$
  \[
    \Im(z) = b
  \]
  \pause
  Os números complexos com parte imaginária nula são os Reais
\end{frame}

%------------------------------------------------------------------------------%
\section{Operações com números complexos}

%------------------------------------------------------------------------------%
\begin{frame}{Igualdade entre números complexos}

  \( z_1 = a + bi \) \; é \emph{igual} a \;
  \( z_2 = p + qi \)

  \pause
  \bigskip
  se, e somente se,
  \[
     a = p
     \qquad\text{e}\qquad
     b = q
  \]
\end{frame}

%------------------------------------------------------------------------------%
\begin{frame}{Soma e subtração de números complexos}
  \onslide<+->{\emph{Soma}}
  \[
    \onslide<+->{z_1 + z_2}
    \onslide<+->{\;=\; (a+b\emph{i}) + (p+q\emph{i})}
    \onslide<+->{\;=\; a + b\emph{i} + p + q\emph{i}}
    \onslide<+->{\;=\; (a+p) + (b+q) \emph{i}}
  \]

  \onslide<+->{\emph{Subtração}}
  \[
    \onslide<+->{z_1 - z_2}
    \onslide<+->{\;=\; (a+b\emph{i}) - (p+q\emph{i})}
    \onslide<+->{\;=\; a+b\emph{i} -p -q\emph{i}}
    \onslide<+->{\;=\; (a-p) + (b-q) \emph{i}}
  \]
\end{frame}

%------------------------------------------------------------------------------%
\begin{frame}{Produto de números complexos}
  \onslide<+->{\emph{Produto}}
  \begin{align*}
    \onslide<+->{z_1z_2}
    \onslide<+->{
      & \;=\; (a+b\emph{i})(p+q\emph{i})} \\
    \onslide<+->{
      & \;=\; a         (p+q\emph{i})
        \;+\; b\emph{i} (p+q\emph{i}) } \\
    \onslide<+->{
      & \;=\; ap + aq\emph{i} + bp\emph{i} + bq\emph{i}^2 } \\
    \onslide<+->{
      & \;=\; ap + aq\emph{i} + bp\emph{i} - bq } \\
    \onslide<+->{
      & \;=\; (ap-bq) \;+\; (aq+bp) \emph{i} }
  \end{align*}
\end{frame}

%------------------------------------------------------------------------------%
\begin{frame}{Conjugado}
  O \emph{conjugado} do número complexo
  {\Large
  \[
    z = a + bi
  \]}
  \pause
  é o número
  {\Large
  \[
    \bar{z} = a - bi
  \]}
\end{frame}

%------------------------------------------------------------------------------%
\begin{frame}{Propriedades dos conjugados}
  \large
  \begin{enumerate}[<+->]
  \setlength\itemsep{1.5em}

  \item \;
  \(
    \bar{\bar{z}} \onslide<+->{= z}
  \)

  \item \;
  se \(\bar{z}=z\) \onslide<+->{então $z$ é real}

  \item \;
  \(
    z\bar{z}
    \onslide<+->{= (a+bi)(a-bi)}
    \onslide<+->{= a^2 -abi +abi - b^2 i^2}
  \) \\[3mm]
  \(
    \hspace{21pt}
    \onslide<+->{= a^2 + b^2}
  \)
  \end{enumerate}
\end{frame}

%------------------------------------------------------------------------------%
\begin{frame}{Divisão de Números Complexos}
  \onslide<+->{Dividir \(z_1=a+bi\) por \(z_2=p+qi\)}
  \begin{align*}
    \onslide<+->{\frac{z_1}{z_2}}
    \onslide<+->{& = \frac{a+bi}{p+qi}}
    \onslide<+->{\;=\; \mathcolor{red}{?}}
  \end{align*}

  \begin{align*}
    \onslide<+->{\frac{z_1}{z_2}}
    \onslide<+->{& = \frac{z_1\,\mathcolor{blue}{\bar{z}_2}}{z_2\,\mathcolor{blue}{\bar{z}_2}}}\\[2mm]
    \onslide<+->{& = \frac{(a+bi)(p-qi)}{(p+qi)(p-qi)}}
    \onslide<+->{  = \frac{ap - aqi + bpi -bqi^2}{p^2+q^2}}
    \onslide<+->{  = \frac{(ap+bq) + (bp-aq)i}{p^2+q^2}} \\[2mm]
    \onslide<+->{& = \left(\frac{ap+bq}{p^2+q^2}\right) + \left(\frac{bp-aq}{p^2+q^2}\right)i}
  \end{align*}
\end{frame}

%------------------------------------------------------------------------------%
\begin{frame}{Corpo complexo}
  A soma e o produto nos complexos tem as mesmas propriedades dos reais

  \pause
  \bigskip
  Por isso o conjunto dos complexos é chamado \emph{corpo}

  \pause
  \bigskip
  Exemplos de corpos: racionais, reais, complexos
\end{frame}

%------------------------------------------------------------------------------%
\include{L-01-Numeros_Complexos-ex2}

%------------------------------------------------------------------------------%
\section{Lista Mínima}

%------------------------------------------------------------------------------%
\begin{frame}{Lista Mínima}
  Atenção: A prova é baseada no livro, não nas apresentações
\end{frame}

%------------------------------------------------------------------------------%
\end{document}
%------------------------------------------------------------------------------%
